\documentclass[12pt]{article}
\usepackage[T1]{fontenc}
\usepackage[utf8]{inputenc}
\usepackage{lmodern}
\usepackage{textcomp}
\usepackage{enumitem}
\usepackage{geometry}
\geometry{margin=1in}
\begin{document}
\begin{center}
Answer Key - Religious Studies - K-12 Level Assessment 17 \
\small (Answers are ordered exactly as in the question paper.)
\end{center}

\section*{Multiple Choice}
\begin{enumerate}[label=\arabic*.]
\item \textbf{Answer:} B
\\ \textit{Options:} A) They worshipped exclusively outside of the temple; B) They believed that nakedness was an elemental expression of non-attachment; C) They included women within the ranks of the leadership; D) They believed that the omniscient Jina must eat and sleep
\\ \textit{Note:} The belief that nakedness represents an elemental expression of non-attachment is a fundamental tenet of the Digambara sect of Jainism, emphasizing renunciation of material possessions and physical desires to achieve spiritual liberation.
\item \textbf{Answer:} C
\\ \textit{Options:} A) 515 BCE; B) 535 BCE; C) 586 BCE; D) 70 CE
\\ \textit{Note:} The Babylonian captivity, also known as the Exile, began in 586 BCE when the Babylonian Empire conquered Jerusalem, leading to the destruction of the First Temple and the forced relocation of a significant portion of the Jewish population.
\item \textbf{Answer:} B
\\ \textit{Options:} A) Sansar; B) Hukam; C) Mantar; D) Karam
\\ \textit{Note:} Hukam refers to the divine order or command in Sikhism, signifying the overarching principle that governs the universe and human life, making it the correct term for the concept of a higher principle in the religion.
\item \textbf{Answer:} D
\\ \textit{Options:} A) Offerings; B) Prayers; C) Apparitions; D) Idols
\\ \textit{Note:} Murtis are physical representations or embodiments of deities in Hinduism, commonly referred to as idols, which serve as focal points for worship and devotion.
\item \textbf{Answer:} C
\\ \textit{Options:} A) Early Third Century BCE; B) Second and First Century BCE; C) Late Sixth Century BCE; D) Fourth and Third Century BCE
\\ \textit{Note:} The late sixth century BCE marks the period when Greek philosophers, such as Thales and Anaximander, began to move away from traditional anthropomorphic depictions of the divine, seeking more abstract and rational understandings of the universe and the divine's role within it.
\end{enumerate}

\section*{True / False}
\begin{enumerate}[label=\arabic*.]
\setcounter{enumi}{5}
\item \textbf{Answer:} True
\\ \textit{Options:} A) True; B) False
\\ \textit{Note:} In Jaina traditions, ajiva is indeed defined as the non-soul, encompassing all non-living entities, which contrasts with jiva, the term for living souls.
\item \textbf{Answer:} False
\\ \textit{Options:} A) True; B) False
\\ \textit{Note:} Guru Nanak referred to the 'divine word' as 'Shabad' in his teachings, not 'Nam', which is a concept related to the name of God in Sikhism but not specifically synonymous with the divine word.
\item \textbf{Answer:} True
\\ \textit{Options:} A) True; B) False
\\ \textit{Note:} The title Dalai Lama is derived from the Mongolian word "dalai," meaning "ocean," and the Tibetan term "lama," meaning "spiritual teacher," which together signify "ocean of wisdom" in Tibetan culture.
\item \textbf{Answer:} False
\\ \textit{Options:} A) True; B) False
\\ \textit{Note:} Confucius is primarily known for his emphasis on ethics, morality, and social harmony rather than mysticism or the concept of qi.
\item \textbf{Answer:} True
\\ \textit{Options:} A) True; B) False
\\ \textit{Note:} Kannon, also known as Avalokiteshvara in Sanskrit, is revered in Japanese Buddhism as the bodhisattva embodying compassion and is dedicated to helping alleviate the suffering of all beings.
\end{enumerate}

\section*{Long Form Response}
\begin{enumerate}[label=\arabic*.]
\setcounter{enumi}{10}
\item \textbf{Answer:} The principles of Zen Buddhism significantly influence the minimalist style of Haiku poetry by emphasizing simplicity, clarity, and a deep connection with nature. Zen encourages practitioners to focus on the present moment, which is reflected in Haiku's concise structure and vivid imagery that captures fleeting experiences. This relationship between Zen and Haiku illustrates how both art forms prioritize mindfulness and the beauty of simplicity, allowing readers to find profound meaning in just a few words. By understanding this connection, we gain insight into how Zen philosophy shapes artistic expression and invites deeper contemplation of life and nature.
\\ \textit{Note:} correct because it effectively connects the principles of Zen Buddhism, such as simplicity and mindfulness, to the concise and nature-focused structure of Haiku poetry, demonstrating how both art forms reflect a shared appreciation for the present moment and the beauty found in simplicity.
\item \textbf{Answer:} Remembering the Divine Name, or Naam, is a central practice in Sikhism that holds great significance for followers. This practice involves constant meditation on and recitation of God's name, which helps individuals develop a deep connection with the Divine. By focusing on the Divine Name, Sikhs believe they can overcome their ego and worldly distractions, leading to spiritual growth and enlightenment. Ultimately, this remembrance fosters a sense of inner peace and unity with God, which is essential for achieving spiritual liberation, or Mukti, in Sikh belief. Thus, the practice of remembering the Divine Name is vital for Sikhs as it guides them on their spiritual journey towards liberation.
\\ \textit{Note:} correct because it emphasizes that remembering the Divine Name in Sikhism facilitates a deep connection with God, aids in overcoming ego and distractions, and ultimately leads to spiritual growth and liberation.
\end{enumerate}

\vfill
\noindent\textit{End of Answer Key}
\end{document}